%=============================================================================
\appendix

\section{Descripteurs de Fukui testés pour la correction de $E_{\text{relax}}$}
\label{app:fukui}

Le tableau~\ref{tab:fukui_exhaustive} récapitule l'ensemble des descripteurs de Fukui et de CDFT testés en corrélation avec le résidu $E_{\text{relax}} - E_{\text{relax}}^{\text{LVC2}}$. Pour chaque descripteur, la corrélation est évaluée séparément pour la réduction (descripteur $f^+$ ou dual) et l'oxydation (descripteur $f^-$ ou dual). Les descripteurs sont regroupés par catégorie. Seuls les descripteurs intégrés sur les oxygènes atteignent la significativité ($p < 0.05$) en réduction ; aucun descripteur n'est significatif en oxydation.

\begin{table}[H]
\centering
\caption{Corrélations entre les descripteurs de Fukui/CDFT et le résidu LVC2 ($n = 13$ molécules). Les valeurs significatives ($p < 0.05$) sont indiquées en gras.}
\label{tab:fukui_exhaustive}
\footnotesize
\begin{tabular}{llrrrr}
\toprule
 & & \multicolumn{2}{c}{Réduction} & \multicolumn{2}{c}{Oxydation} \\
\cmidrule(lr){3-4} \cmidrule(lr){5-6}
Catégorie & Descripteur & $r$ & $p$ & $r$ & $p$ \\
\midrule
\multirow{3}{*}{Intégrés sur O}
& $f^{\pm}_{\text{sumO}}$ & $\mathbf{-0.67}$ & $\mathbf{0.012}$ & $-0.16$ & $0.59$ \\
& $\Delta f_{\text{sumO}}$ & $\mathbf{-0.69}$ & $\mathbf{0.009}$ & $-0.35$ & $0.24$ \\
& $f^0_{\text{sumO}}$ & $\mathbf{-0.67}$ & $\mathbf{0.012}$ & $-0.16$ & $0.59$ \\
\midrule
\multirow{2}{*}{Intégrés sur C}
& $f^{\pm}_{\text{sumC}}$ & $+0.54$ & $0.058$ & $+0.07$ & $0.82$ \\
& $\Delta f_{\text{sumC}}$ & $+0.54$ & $0.058$ & $+0.36$ & $0.22$ \\
\midrule
\multirow{2}{*}{Intégrés sur H}
& $f^{\pm}_{\text{sumH}}$ & $+0.04$ & $0.88$ & $-0.36$ & $0.23$ \\
& $\Delta f_{\text{sumH}}$ & $+0.04$ & $0.88$ & $-0.36$ & $0.23$ \\
\midrule
\multirow{2}{*}{Atomiques (O)}
& $\max_A f^{\pm}_A$ (O) & $-0.34$ & $0.25$ & $-0.16$ & $0.60$ \\
& $\Delta f_{\max}$ / $\Delta f_{\min}$ & $-0.34$ & $0.25$ & $+0.39$ & $0.19$ \\
\midrule
\multirow{3}{*}{Top-$k$ atomiques}
& Top-2 $\Delta f$ & $-0.41$ & $0.16$ & $+0.38$ & $0.20$ \\
& Top-3 $\Delta f$ & $-0.38$ & $0.20$ & $+0.41$ & $0.17$ \\
& Top-5 $\Delta f$ & $-0.31$ & $0.30$ & $+0.35$ & $0.24$ \\
\midrule
\multirow{3}{*}{CDFT globaux}
& $\mu$ (pot. chimique) & $-0.14$ & $0.66$ & $-0.25$ & $0.40$ \\
& $\eta$ (dureté) & $-0.30$ & $0.32$ & $-0.44$ & $0.13$ \\
& $\omega$ (électrophilicité) & $+0.20$ & $0.52$ & --- & --- \\
\midrule
Référence & $f_{\text{O}}^{\text{LUMO}}$ ($n = 15$) & $\mathbf{-0.68}$ & $\mathbf{0.005}$ & $-0.49$ & $0.061$ \\
\bottomrule
\end{tabular}
\end{table}

\section{Exemples de Fichiers d'Entrée AMS}
%=============================================================================

Les exemples suivants illustrent les fichiers d'entrée générés par PLAMS pour chaque pipeline.

\subsection{DFT Optimisation + Fréquences (Pipeline Conformères, étape 4)}

\begin{lstlisting}
Task GeometryOptimization

System
  Atoms
    C   0.000  0.000  0.000
    ...
  End
  Charge 0
End

Properties
  NormalModes Yes
End

Engine ADF
  Basis
    Type DZP
    Core None
  End
  Relativity
    Level None
  End
  NumericalQuality Good
  XC
    Hybrid PBE0
  End
  Solvation
    Surf Delley
    Solv name=Acetonitrile
    Charged method=CONJ
    C-Mat POT
    SCF VAR ALL
  End
EndEngine
\end{lstlisting}

\subsection{Optimisation Géométrique DZP (Pipeline Redox)}

Cet exemple montre l'optimisation géométrique DZP utilisée pour les états chargés (reduced\_opt, oxidized\_opt). Le neutre n'est pas re-optimisé (réutilisation du RKF conformères) :

\begin{lstlisting}
Task GeometryOptimization

System
  Atoms
    ...
  End
  Charge -1       ! -1 (reduit) ou +1 (oxyde)
End

Properties
  NormalModes Yes
End

Engine ADF
  Basis
    Type DZP
    Core None
  End
  Relativity
    Level None
  End
  Unrestricted Yes
  SpinPolarization 1
  XC
    Hybrid PBE0
  End
  Solvation
    Surf Delley
    Solv name=Acetonitrile
    Charged method=CONJ CORR  ! CORR pour les ions
    C-Mat POT
    SCF VAR ALL
  End
EndEngine
\end{lstlisting}

\subsection{Single-Point TZ2P (Pipeline Redox)}

Les single-points TZ2P sont calculés pour le raffinement énergétique du schéma composite. Pas de modes normaux :

\begin{lstlisting}
Task SinglePoint

System
  Atoms
    ...  ! Geometrie neutre ou ionique optimisee
  End
  Charge -1  ! 0 (neutral_tz2p), -1 (reduced_sp/tz2p), +1 (oxidized_sp/tz2p)
End

Engine ADF
  Basis
    Type TZ2P
    Core None
  End
  Relativity
    Level None
  End
  Unrestricted Yes   ! No pour neutral_tz2p
  SpinPolarization 1 ! 0 pour neutral_tz2p
  XC
    Hybrid PBE0
  End
  Solvation
    Surf Delley
    Solv name=Acetonitrile
    Charged method=CONJ CORR  ! CONJ pour neutral_tz2p
    C-Mat POT
    SCF VAR ALL
  End
EndEngine
\end{lstlisting}

\subsection{Analyse QTAIM/IQA}

\begin{lstlisting}
Task SinglePoint

System
  Atoms
    ...
  End
  Charge 0  ! 0, -1 ou +1 selon l'etat analyse
End

Engine ADF
  Basis
    Type TZP
    Core None
  End
  Relativity
    Level None
  End
  NumericalQuality Good
  XC
    Hybrid PBE0
  End
  QTAIM
    Enabled Yes
    AnalysisLevel Full
    Source Yes
  End
  IQA
    Enabled Yes
  End
EndEngine
\end{lstlisting}

